\chapter{Wiring the reader}
\label{ch:wiring}

\chapterbrief
  {Connect the RC522 and prove it answers.}
  {Separating ``is it the wiring?'' from ``is it the code?'' --- permanently.}
  {\code{WIRING OK. The RC522 is alive and talking.}}

This is the last chapter of Volume Zero. When its checkpoint passes, you have a
complete working environment and you can start writing code.

\section{Power off first}

\begin{shell}
ssh brewer 'sudo shutdown -h now'
\end{shell}

Wait for the green LED to stop flickering, then unplug the power. Wiring a
powered Pi is how you destroy a Pi.

\section{The seven wires}

\begin{center}\small
\begin{tabular}{@{}llll@{}}
\toprule
\thead{RC522} & \thead{Pi pin} & \thead{Pi name} & \thead{Note} \\
\midrule
\reg{3.3V} & 1  & 3V3           & \textbf{Never 5\,V} \\
\reg{RST}  & 22 & BCM 25        & reset \\
\reg{GND}  & 6  & GND           & \\
\reg{MISO} & 21 & BCM 9         & chip $\rightarrow$ Pi \\
\reg{MOSI} & 19 & BCM 10        & Pi $\rightarrow$ chip \\
\reg{SCK}  & 23 & BCM 11        & clock \\
\reg{SDA}  & 24 & BCM 8 (CE0)   & chip select \\
\reg{IRQ}  & --- & ---          & leave unconnected \\
\bottomrule
\end{tabular}
\end{center}

\begin{center}
\begin{tikzpicture}[font=\scriptsize\sffamily, x=1cm, y=1cm]
  % Pi header: two columns of pins, numbered as on the board
  \draw[brewernavy, line width=0.8pt, rounded corners=2pt] (0,-0.3) rectangle (1.5,6.3);
  \node[color=brewernavy, font=\scriptsize\sffamily\bfseries] at (0.75,6.6) {Pi header};
  \foreach \y/\l/\r in {6/1/2, 5.5/3/4, 5/5/6, 4.5/7/8, 4/9/10,
                        3.5/11/12, 3/13/14, 2.5/15/16, 2/17/18,
                        1.5/19/20, 1/21/22, 0.5/23/24, 0/25/26}{
     \fill[brewergrey!45] (0.45,\y) circle (1.5pt);
     \fill[brewergrey!45] (1.05,\y) circle (1.5pt);
     \node[color=brewergrey!70, font=\tiny] at (0.2,\y) {\l};
     \node[color=brewergrey!70, font=\tiny] at (1.3,\y) {\r};
  }
  % highlight the pins we use
  \foreach \y/\x/\c in {6/0.45/brewerred, 5/1.05/brewerink, 1.5/0.45/brewerblue,
                        1/1.05/breweramber, 1/0.45/brewergreen,
                        0.5/0.45/brewerblue, 0.5/1.05/brewergreen}{
     \fill[\c] (\x,\y) circle (2.4pt);
  }
  % RC522
  \draw[breweramber!80!black, line width=0.8pt, rounded corners=2pt] (5.6,0.4) rectangle (7.4,5.6);
  \node[color=breweramber!80!black, font=\scriptsize\sffamily\bfseries] at (6.5,5.9) {RC522};
  \foreach \y/\lbl/\c in {5.2/3.3V/brewerred, 4.6/RST/breweramber,
                          4.0/GND/brewerink, 3.4/MISO/brewergreen,
                          2.8/MOSI/brewerblue, 2.2/SCK/brewerblue,
                          1.6/SDA/brewergreen, 1.0/IRQ/brewergrey}{
     \fill[\c] (5.75,\y) circle (2.2pt);
     \node[color=brewergrey, anchor=west, font=\tiny] at (5.95,\y) {\lbl};
  }
  % wires (schematic, not geometric)
  \draw[brewerred,   line width=0.8pt] (0.45,6)   -- (4.2,6)   -- (4.2,5.2) -- (5.75,5.2);
  \draw[breweramber, line width=0.8pt] (1.05,1)   -- (3.4,1)   -- (3.4,4.6) -- (5.75,4.6);
  \draw[brewerink,   line width=0.8pt] (1.05,5)   -- (4.6,5)   -- (4.6,4.0) -- (5.75,4.0);
  \draw[brewergreen, line width=0.8pt] (0.45,1)   -- (3.0,1)   -- (3.0,3.4) -- (5.75,3.4);
  \draw[brewerblue,  line width=0.8pt] (0.45,1.5) -- (2.6,1.5) -- (2.6,2.8) -- (5.75,2.8);
  \draw[brewerblue,  line width=0.8pt] (0.45,0.5) -- (2.2,0.5) -- (2.2,2.2) -- (5.75,2.2);
  \draw[brewergreen, line width=0.8pt] (1.05,0.5) -- (1.8,0.5) -- (1.8,1.6) -- (5.75,1.6);
  \node[color=brewergrey, font=\tiny, anchor=west] at (2.0,-0.55)
       {pin 1 is the corner nearest the SD card slot};
\end{tikzpicture}
\end{center}

\begin{trap}{3.3\,V, not 5\,V}
Pin 1 is 3.3\,V. Pin 2 --- \emph{directly beside it} --- is 5\,V, and it will
destroy the RC522.

Count carefully. Pin 1 is the corner nearest the SD card slot; odd numbers run
down one side, even down the other. Getting this wrong is silent, immediate and
permanent.
\end{trap}

\begin{trap}{\reg{SDA} is not I\textsuperscript{2}C here}
The RC522 board labels its chip-select pin \reg{SDA}, which is an
I\textsuperscript{2}C name. The module supports several buses and the silkscreen
reflects a different one.

For SPI, \reg{SDA} goes to \textbf{CE0, Pi pin 24}. This is the single most
common mis-wiring, and its symptom is \code{VersionReg = 0x00}.
\end{trap}

\section{Check it twice}

Before powering on, verify each wire by touching both ends, one at a time. Seven
wires, seven checks. Two minutes now saves an hour later --- and unlike code,
you cannot \code{printf} a wire.

Then power up and wait a minute.

\section{Probe it}

\begin{shell}
# on your laptop
make deploy-full
make probe
\end{shell}

\begin{verifybox}
\begin{trace}
=== RC522 wiring probe ===
Opening SPI bus 0, chip-select 0, 1 MHz, mode 0...
VersionReg   (0x37) = 0x92
TxControlReg (0x14) = 0x80

  -> WIRING OK. The RC522 is alive and talking.
     You are ready for Chapter 1 of the main book.
\end{trace}
\end{verifybox}

\hex{0x92} is the MFRC522 v2.0 identifier (\hex{0x91} is v1.0 --- also fine).
The chip is powered, on the right pins, and speaking SPI at the right clock and
mode.

\section{If it does not say that}

The probe diagnoses itself. Its three failure messages, and what they mean:

\begin{center}\small
\begin{tabular}{@{}L{0.16\linewidth}L{0.32\linewidth}L{0.44\linewidth}@{}}
\toprule
\thead{Reading} & \thead{Meaning} & \thead{Check, in order} \\
\midrule
\hex{0x00} & MISO stuck low --- the chip is not driving the line &
  Power (3.3\,V and GND); then \reg{SDA} $\rightarrow$ pin 24; then \reg{MISO}
  $\rightarrow$ pin 21 \\
\hex{0xFF} & MISO floating --- nothing connected &
  \reg{MISO} $\rightarrow$ pin 21. If connected, suspect a dead module \\
varies & noise & \textbf{Shorter wires.} Run the probe three times: if the value
  \emph{changes}, it is definitely wiring \\
\bottomrule
\end{tabular}
\end{center}

And if it fails to open SPI at all:

\begin{center}\small
\begin{tabular}{@{}L{0.40\linewidth}L{0.52\linewidth}@{}}
\toprule
\thead{Error} & \thead{Fix} \\
\midrule
\code{No such file or directory} & SPI not enabled. \code{sudo raspi-config},
  Interface Options, then \textbf{reboot} \\
\code{Permission denied} & Not in the \code{spi} group --- and remember the
  change needs a fresh login \\
\code{No gpio chip found} & You ran it on your laptop, not the Pi \\
\bottomrule
\end{tabular}
\end{center}

\begin{concept}{Why this probe exists at all}
It answers one question and nothing else: \emph{is the hardware wired correctly?}

From Volume One, Chapter 1 onwards you will be writing your own driver, and its
bugs will look exactly like wiring faults. Having a known-good, disposable
diagnostic means you never have to wonder which one you are looking at.

Keep it. Run it whenever anything hardware-shaped goes wrong. It is deliberately
ugly --- one flat method, no abstractions --- because a diagnostic you have to
debug is worthless.
\end{concept}

\section{You are done}

\begin{checkpoint}{Volume Zero complete}
\begin{shell}
ssh brewer 'hostname'          # one word, no password
cd ~/midnight-brewer && mvn test    # Tests run: 3
make deploy && make run        # runs on the Pi, aarch64
make probe                     # WIRING OK
git log --oneline              # at least one commit
git push                       # backed up
\end{shell}

Six commands, six green results. You have:

\begin{itemize}
  \item a headless Pi reachable by name, with key-only SSH;
  \item Java 21 and Maven on both machines;
  \item a project that builds and tests in seconds;
  \item a five-second deploy loop to real hardware;
  \item your work in git and on GitHub;
  \item an NFC reader answering.
\end{itemize}

None of that was specific to NFC. It is the setup for any embedded Linux project
you will ever do.
\end{checkpoint}

\vspace{1em}

Now open \textbf{Volume One} --- \emph{Building a Contactless Payment Terminal
From Scratch} --- and turn to Chapter 1. You will write an interface, build a
fake SPI bus, and test a driver with the reader unplugged.

\begin{flushright}\itshape
The environment is no longer something that can be wrong.
\end{flushright}
